\documentclass[11pt]{article}

\usepackage[margin=1in]{geometry}
\usepackage{amsmath,amssymb,amsthm}
\usepackage{booktabs}
\usepackage{microtype}

\newtheorem{theorem}{Theorem}
\newtheorem{remark}{Remark}

\newcommand{\oracle}{\widehat{O}}
\newcommand{\eps}{\varepsilon}

\title{Constraint-first string theory:\\ the flux landscape's only exact structure is its constraint}
\author{Rob Mitchell}
\date{}

\begin{document}
\maketitle

\begin{abstract}
We measure, on real supersymmetric flux vacua, which low-energy observables carry
information about the flux that produced them. Every flux-determined observable is
classified by a joint (multi-coordinate quadratic) learnability test, gated by a
positive control that must reproduce the one exactly-known quantity. The result is a
\emph{structure map}: the string coupling, the gauge couplings, the mass-spectrum
\emph{shape}, the axions, the dilaton, and the cosmological constant are all
statistically uniform in flux ($r < 0.25$); $W_0$ and the saxions carry weak
structure ($r \approx 0.3$--$0.6$); and exactly one quantity is \emph{exactly}
structured --- the D3 tadpole, the symplectic pairing $N_{\rm flux} = F^\top J H$
($r = 1.0$). No observable carries the strong ($r > 0.9$) \emph{geometric
fingerprint} that would pin the compactification. We conclude: the landscape's only
exact structure is what it \emph{forbids}, not what it \emph{selects}; everything it
produces is noise plus a weak scale. This is the constraint-first program,
established by measurement rather than asserted.
\end{abstract}

\section{Introduction}

The string landscape is $\sim 10^{500}$ vacua and no experiment that tells most of
them apart. Two programs have tried to make it navigable. The \emph{search} program
(and its modern machine-learning descendants) asks: can we \emph{learn} the map from
flux to observables, and so scan our way to our universe? The \emph{constraint}
program (swampland, bootstrap) asks instead: which consistency conditions
\emph{eliminate} vacua? This paper measures the premise that distinguishes them:
whether any observable carries a \emph{fingerprint} --- strong information about the
flux --- that a search could exploit.

The answer, measured on $174$ supersymmetric IIB flux vacua, is that no observable
does. The only exact structure is the constraint itself.

\section{Method}

\paragraph{Sample.} JAXVacua~\cite{jaxvacua} on a bundled $h^{1,2}=2$ flux EFT, the
natural ISD measure (${\tt isd\_mode}={\tt ISD-}$, Gaussian flux prior),
degree-$0$ spectrum. $174$ supersymmetric vacua (selected by $DW=0$,
$\operatorname{Im}\tau>0$) from $1318$ critical points.

\paragraph{The joint test and the Rule of Three.} For each observable $y$ we compute
the cross-validated ridge regression of $y$ on the degree-$2$ monomial basis of the
$12$ flux coordinates ($91$ features), reporting the Pearson $r$ between held-out
predictions and truth. A single-coordinate mutual-information test is \emph{blind} to
multi-coordinate structure and is used only as a first screen; the claim follows the
\emph{Rule of Three}: single-coordinate MI, joint ridge-$r$, and a permutation null
must \emph{agree}. A disagreement is unresolved, never resolved by one test.

\paragraph{The control.} The D3 tadpole is \emph{known} to be the symplectic pairing
$N_{\rm flux} = F^\top J H$ with $J = \bigl(\begin{smallmatrix}0&I\\-I&0\end{smallmatrix}\bigr)$,
a clean bilinear form of the flux. The joint test must return $r = 1.0$ for it, at
machine precision. This control is what makes every other reading trustworthy: a
test that cannot see the tadpole's structure cannot be trusted to see any.

\section{Results: the structure map}

\begin{center}
\begin{tabular}{lcc}
\toprule
observable & wide-flux $r$ & verdict \\
\midrule
tadpole $N_{\rm flux}$ & \textbf{1.00} & \textbf{exact} (the constraint) \\
$W_0$ & 0.62 & weakly learnable (scale via tadpole) \\
saxions (moduli sizes) & 0.30--0.42 & weakly learnable (stabilization) \\
string coupling $g_s$ & 0.12 & uniform \\
gauge couplings & 0.02--0.12 & uniform \\
mass-spectrum shape (ratios) & $<$0.25 & uniform \\
axions (moduli phases) & 0.15 & uniform \\
axio-dilaton $\operatorname{Re}\tau$ & 0.13 & uniform \\
AdS potential $V$ & 0.00 & uniform \\
\bottomrule
\end{tabular}
\end{center}

The pattern is stark. The low-energy observables are a mostly-unstructured cloud:
the string coupling, the gauge couplings, the mass-spectrum \emph{shape}, the
axions, the dilaton, and the cosmological constant are statistically uniform in
flux. Two things carry weak structure --- $W_0$, whose \emph{scale} tracks the
tadpole, and the saxions, which flux stabilizes. Exactly one thing is \emph{exactly}
structured: the tadpole, the constraint.

\begin{remark}
\textbf{No geometric fingerprint exists.} Across every observable measured, none
carries strong ($r>0.9$) information about the flux. The fingerprint that a search
program would need --- an observable that pins the compactification --- is absent.
The only exact structure is what the landscape \emph{forbids}.
\end{remark}

\section{The $W_0$ correction and the search}

A first pass reported $W_0$ as ``zero-information,'' on the strength of a
single-coordinate correlation; that is retracted here. The joint test gives
$r = 0.62$ on the natural measure ($r=0.10$ only when flux is artificially bounded).
$W_0$ is weakly learnable, through its scale.

This does not rescue search. On the sample, $\log|W_0|$ has standard deviation
$0.63$ nats, so a model with $r=0.62$ has prediction error $\sigma_{\rm pred}=0.50$
nats. The KKLT tuning target $\log|W_0|=-9.2$ sits $9.9$ nats --- about
$20\,\sigma_{\rm pred}$ --- below the mean. A model fit on the data cannot guide a
search to a point twenty prediction-errors away; that is a meaningless
extrapolation, not a gradient.

\begin{remark}
\textbf{Search over $W_0$ is hopeless by rarity plus extrapolation}, not by
zero-information. $W_0$ is log-uniform (CDF $R^2=0.87$--$0.92$), so small $W_0$ is
exponentially rare; and the target is $20\,\sigma_{\rm pred}$ outside the observed
range. The conclusion survives; the argument is now measured.
\end{remark}

\section{Conclusion}

The constraint-first program is established by measurement:

> \textbf{The flux landscape's only exact structure is its constraint.} Everything
> the landscape \emph{produces} --- couplings, masses, phases, the cosmological
> constant --- is noise plus a weak scale; the one thing that is exactly determined
> is the D3 tadpole, the quantity that \emph{rules vacua out}. No search can find our
> vacuum in observables, because there is nothing structured to find; the only
> navigable move is elimination by verified constraint.

\paragraph{Scope, named.} One model ($h^{1,2}=2$), degree-$0$ spectrum, $174$ vacua,
flux coordinates only, the IIB flux corner. The instanton-corrected (K\"ahler-moduli)
masses remain computationally infeasible. Each of these is a named limit, not a
hidden one; the structure map is a measurement to be reproduced, and the Rule of
Three is the instrument that keeps it honest.

\begin{thebibliography}{9}
\bibitem{jaxvacua} A.~Dubey, S.~Krippendorf, A.~Schachner, ``JAXVacua'' (arXiv:2306.06160).
\bibitem{denef} F.~Denef and M.~R.~Douglas, ``Distributions of flux vacua,'' JHEP \textbf{0405} (2004) 072.
\bibitem{valiant} L.~G.~Valiant, ``A theory of the learnable,'' CACM \textbf{27} (1984) 1134.
\bibitem{valiant2009} L.~G.~Valiant, ``Evolvability,'' J. ACM \textbf{56} (2009) 3.
\end{thebibliography}

\end{document}
